\section{The unramified Hermitian mass ratio}
\label{sec:mass-formula}

\status{MF1--MF20 reconstruct the actual local densities, the local
orbit identity, the singular-series average, and the rank induction
behind \texttt{SiegelMassRatioCM}. The only named external theorem in
the global representation step is the classical Hasse--Minkowski
local-to-global theorem for quadratic forms over number fields.
Its statement and exact use are recorded in MF4; its proof is
not claimed here. No mass formula, Tamagawa-number theorem,
representation-average formula, or quantitative prime estimate is
assumed. The mass premise retained in QW3 and QW9 is reduced to this
explicitly identified classical theorem. The separate geometric
volume and compatible arithmetic-data premises remain.}

Write \(E\) for the source field \(K\), \(F=E^{\bar{\ }}\) for its
fixed field, \(f=[F:\mathbb Q]\), and
\(v=|\operatorname{disc}E|^{1/(2f)}/2\).
All Hermitian forms are conjugate linear in the first variable.
An integral lattice is a full \(\mathcal O_E\)-lattice on which
all pairings are integral; it need not be free globally. The mass is
the literal sum \(\sum_{[L']}|\operatorname{Aut}(L')|^{-1}\) over the
genus in the given rational Hermitian space, as in IG9.
At a finite place \(\mathfrak p\) of \(F\), put
\[
 R=\widehat{\mathcal O}_{F,\mathfrak p},\quad
 B=\mathcal O_E\otimes_{\mathcal O_F}R,\quad
 q=N\mathfrak p,\quad \eta_{\mathfrak p}=
 \begin{cases}1,&B=R\times R,\\-1,&B/R\text{ is unramified quadratic}.
 \end{cases}
\]
Here \(\pi\) is an actual local uniformizer and
\(\operatorname{ord}_{\mathfrak p}\pi=1\).
The source's congruences and localizations are over the fixed-field
prime \(\mathfrak p\); a split pair of primes of \(E\) is treated
together. No rational prime is identified with \(\pi\).

\paragraph{MF1: the fixed field, the different, and the normalizations.}
The literal \texttt{IsCMUnram} hypothesis gives complex conjugation
under every embedding and, at every prime \(P\) of \(E\), an integral
\(\theta\) with \(\theta-\bar\theta\notin P\).
Choose a nonzero \(w\) with \(\bar w=-w\). Such a \(w\) exists because
otherwise every embedded element of \(E\) would be real.
Every \(z\in E\) has the unique decomposition
\[
 z=\frac{z+\bar z}{2}
       +w\frac{z-\bar z}{2w},\qquad E=F\oplus Fw.       \tag{MF.1}
\]
Thus \(E/F\) is quadratic, \(F\) is totally real, and \(E\) is totally
imaginary. At a fixed finite prime the element
\(\theta-\bar\theta\) is a unit at both conjugate primes.
The coefficient formulas in DR1, applied to \(1,\theta\), show that
the integral quadratic algebra over \(R\) has this basis. Its
discriminant \((\theta-\bar\theta)^2\) is a unit. Hence it is either
split or unramified quadratic, also over primes above \(2\).
Its trace pairing is perfect and its relative different is the unit
ideal. Taking determinants of the successive trace pairings gives
\[
 |\operatorname{disc}E|=(\operatorname{disc}F)^2,\quad
 \operatorname{disc}F=(2v)^f,\quad
 \operatorname{covol}(\mathcal O_E)=v^f,\quad
 \operatorname{covol}(\mathcal O_F)=\sqrt{(2v)^f}.       \tag{MF.2}
\]
The complex covolume uses ordinary Lebesgue measure on one copy of
\(\mathbb C\) per conjugate pair: it is
\(2^{-f}\sqrt{|\operatorname{disc}E|}\). The real covolume is
\(\sqrt{\operatorname{disc}F}\).

Set
\[
 A_m(\mathfrak p)=\prod_{j=1}^{m}(1-\eta_{\mathfrak p}^{\,j}q^{-j}),
 \quad e_m(\mathfrak p)=1-\eta_{\mathfrak p}^{\,m}q^{-m},
 \quad L_m=\prod_{\mathfrak p}e_m(\mathfrak p)^{-1}.     \tag{MF.3}
\]
For \(m\ge2\), this product is positive and absolutely convergent
after taking logarithms, by QW1. It is exactly
\(\texttt{lValEta}(E,m)\): an inert prime of \(F\) has one prime of
\(E\) of norm \(q^2\), while a split prime gives two factors with
exponent \(-1/2\). Their product is (MF.3), with no residual
square-root factor.

\paragraph{MF2: local Hermitian diagonalization, including dyadic places.}
The norm \(B^\times\to R^\times\) is onto.
In the nonsplit case first solve the norm equation in
\(\mathbb F_{q^2}^\times\). If a unit has the desired norm modulo
\(\pi^j\), multiply it by \(1+\pi^j t\); the norm changes modulo
\(\pi^{j+1}\) by \(1+\pi^j\operatorname{Tr}(t)\).
The residue trace is onto, even in characteristic \(2\), so successive
corrections converge. In the split case the norm is \((a,b)\mapsto ab\).
This is the norm-lifting argument of DR1; it never divides by \(2\).

Every integral nondegenerate Hermitian lattice over nonsplit \(B\)
has an orthogonal basis with Gram matrix
\(\operatorname{diag}(\pi^{a_1},\ldots,\pi^{a_m})\), \(a_i\ge0\).
Choose a pairing of least valuation \(a\).
If no diagonal entry has that valuation, replacing a vector by
\(x+t y\) produces one: the leading cross term is a nonzero residue
trace for a suitable \(t\). Projection along this vector has integral
coefficients, since all pairings are divisible by \(\pi^a\), so
induction splits off the line. Unit-norm surjectivity removes each
diagonal unit. For split \(B\), the two components of a lattice are
free \(R\)-modules paired by a matrix; Smith reduction gives the same
description, with independent bases in the two components.
The multiset of \(a_i\) is intrinsic, as the elementary divisors of
the map to the dual lattice.

Over the fraction algebra, a nonsplit form has a basis with all
exponents \(0\), except possibly one exponent \(1\).
First reduce exponents modulo \(2\). Two entries \(\pi,\pi\)
can be replaced by \(1,1\): choose a unit \(t\) with
\(N(t)=\pi-1\); then
\(\pi N(\pi^{-1})+\pi N(\pi^{-1}t)=1\).
Split off this norm-one vector. The complementary coefficient has
even valuation, hence is a norm and can also be normalized to \(1\).
In the split case every fraction-field form is standard.
Consequently an integral lattice of determinant valuation at most
one is determined up to integral isometry by its rank and determinant
valuation. If a fixed vector has norm one, keep it as the first
basis vector and choose such a minimal lattice in its orthogonal
complement. This proves the local construction used in
\texttt{exists\_lattice\_vd\_le\_one}.

\paragraph{MF3: exact norm orbits and rational local approximation.}
For vectors \(x,y\) of the same nonzero norm, put \(u=x-y\).
If \(c=h(u,x)\ne0\), the map
\[
             z\longmapsto z-u\,c^{-1}h(u,z)            \tag{MF.4}
\]
is an isometry taking \(x\) to \(y\). Indeed
\(c+\bar c=h(u,u)\), which cancels the three extra terms on
expanding the pairing. If \(c=0\), pass through \(-y\);
the corresponding denominator is \(2h(x,x)\ne0\), and multiplication
by \(-1\) is an isometry. Thus the unitary group is transitive on
each nonzero-norm fiber. Over a positive definite global space,
the first denominator can vanish only when \(x=y\).

The rational unitary group is dense in its group over any one
completion, and simultaneously over any prescribed finite set.
On skew-adjoint matrices \(Z^*=-Z\), where here the star is the
adjoint for the fixed Hermitian form, the Cayley transform
\[
             Z\longmapsto (1-Z)(1+Z)^{-1}              \tag{MF.5}
\]
is rational, with inverse \((1-g)(1+g)^{-1}\) when
\(\det(1+g)\ne0\). The skew-adjoint matrices form an \(F\)-vector
space, so ordinary weak approximation applies to their coordinates.
To enter this chart near an arbitrary \(g\), multiply \(g\) by a
norm-one scalar arbitrarily close to \(1\) avoiding the finitely
many forbidden eigenvalues. Such scalars are supplied by the
one-dimensional Cayley transform on \(Fw\); weak approximation on
\(Fw\) chooses one rational scalar for all the specified places.
Apply weak approximation
to the resulting chart coordinates and then undo the scalar.
This proves the assertion without strong approximation.

Lattice stabilizers are open. A completed isometry carrying one
rational localized lattice to another can therefore be replaced
by a rational isometry with exactly the same lattice action.
The same argument applies to the orthogonal complement in (MF.4).
Thus the completed orbits below are exactly the rational local
orbits in \texttt{fibLoc}, \texttt{IsLocReps}, and
\texttt{SameGenus}. Those definitions do not merely ask for completed
isometries.

\paragraph{MF4: the binary representation theorem and its classical input.}
We use the following named classical theorem:
\emph{a nondegenerate quadratic form over a number field has a nonzero
isotropic vector if it does over every completion}.
This is Hasse--Minkowski in its isotropy formulation; the
primary references are Kim, \emph{Hasse--Minkowski Theorem},
Theorem~1.1, and Hanke, \emph{Quadratic Forms and Automorphic Forms},
Theorem~1.9.1 and Remark~1.9.2 (printed p.~20).\footnote{Author-written
texts: \href{https://www.csun.edu/~sungjin/HasseMinkowski.pdf}{Kim's
statement} and \href{https://swc-math.github.io/aws/2009/09HankeNotes.pdf}{Hanke's
expanded Arizona Winter School notes}. The report distinguishes
their stated general theorem from proofs specialized to \(\mathbb Q\).}
Its global proof is not included here.

Let \(a,b,n\in F\) be totally positive. At a finite place, the binary
Hermitian form \(aN(s)+bN(t)\) represents every element of
\(F_{\mathfrak p}\).
For split \(E_{\mathfrak p}\), one norm already represents every
element. In the unramified field case, remove norm factors from
\(a,b\), leaving exponents \(0\) or \(1\).
If their parities differ, one coefficient has the target's valuation
parity, and unit-norm surjectivity gives a representation with that
coordinate. If their parities agree, the two-dimensional norm sum
is isotropic because \(-1\) is a unit norm, and represents every
target: after absorbing that norm, use
\(N((c+1)/2)-N((c-1)/2)=c\).
This is a fraction-field identity, with no integral division by \(2\).

Choose \(w\) as in (MF.1), with \(\delta=w^2\in F\).
At every real embedding \(\delta<0\).
The five-variable form
\[
 a(X_1^2-\delta X_2^2)+b(X_3^2-\delta X_4^2)-nX_5^2   \tag{MF.6}
\]
is locally isotropic at finite places by the explicit representations,
and at real places because its signature is \((4,1)\).
Hasse--Minkowski gives a nonzero global zero. Its last coordinate
cannot vanish: positivity at a real embedding would then force the
other four coordinates to vanish. Set
\(s=(X_1+wX_2)/X_5,\ t=(X_3+wX_4)/X_5\).
Then \(a\,s\bar s+b\,t\bar t=n\).
This proves precisely \texttt{HermRepUnram}.
Gram--Schmidt now shows that every totally positive definite
Hermitian space of rank at least two represents every totally
positive element of \(F\). Only this global step and its descendants
use the named theorem.

\paragraph{MF5: the matrix-density measure and the local ratio law.}
For a local integral lattice \(L\) of rank \(m\), let \(H\) be its
Gram matrix and \(d=\operatorname{ord}_{\mathfrak p}\det H\).
The source's \(d=\texttt{vd}\) is exactly this integer: determinants
of all Gram matrices of \(m\)-tuples generate \((\det H)\).
Write
\[
 \alpha_L=\lim_{k\to\infty}
 \frac{\#\{X\bmod\pi^k:X^*HX\equiv H\bmod\pi^k\}}{q^{km^2}},
 \qquad \nu(L)=q^{-md}\alpha_L.                       \tag{MF.7}
\]
These are the literal \texttt{akCount}/\texttt{alphaL}
normalizations; the numerator initially counts endomorphisms.
For \(k>d\), determinants show that all such endomorphisms are
invertible. In the split algebra, their two nonnegative determinant
valuations have sum zero.

On the full matrix space give the integral matrix lattice volume
one; do the same on the \(F_{\mathfrak p}\)-space of Hermitian
matrices. The map \(X\mapsto X^*HX-H\) has surjective derivative at
every unitary \(X\). At \(1\) it is \(Z\mapsto Z^*H+HZ\).
To solve it, write a Hermitian target as \(Y+Y^*\):
use upper-triangular entries and an integral trace-one element on
the diagonal, and then multiply by \(H^{-1}\).
This exhibits bounded denominators, including at \(2\).
On a sufficiently small ball a linear coordinate change puts the
map into projection form. Successive corrections prove this:
after removing the linear error the quadratic error has strictly
larger valuation, and the corrections converge.
The quotient of ambient measure by target measure is then the
measure on the fiber. Finitely many such balls cover the compact
unitary stabilizer; compactness puts all sufficiently accurate
approximate solutions in their union. Hence (MF.7) stabilizes,
equals its fiber measure, and is positive.

If the basis changes by \(C\), conjugation on the full matrix space
has Jacobian one, whereas \(Y\mapsto C^*YC\) on Hermitian matrices
has Jacobian \(|N_{E/F}\det C|^m\).
For diagonal \(C\), the diagonal and off-diagonal coordinates give
this directly; elementary matrices have determinant one, so
elimination proves the general case.
The determinant valuation changes by
\(\operatorname{ord}_{\mathfrak p}N_{E/F}\det C\).
Therefore \(q^{-md}\) times the fiber measure is basis independent.
Unitary left and right translations preserve it since
\(N_{E/F}\det g=1\). We have constructed a Haar measure satisfying
\[
                \operatorname{vol}(K_L)=\nu(L).       \tag{MF.8}
\]
For two lattices, take the volume of the common finite-index
stabilizer to obtain
\[
 \nu(A)\,\operatorname{orb}(B,A)
          =\nu(B)\,\operatorname{orb}(A,B),             \tag{MF.9}
\]
where \(\operatorname{orb}(A,B)=[K_A:K_A\cap K_B]\).
This is exactly \texttt{LocalDensityRatio} after multiplication by
\(q^{m(d_A+d_B)}\), with the source's orbit orientation.

\paragraph{MF6: the unimodular factors and the full good-prime norm density.}
For the standard form over the residue algebra, the number of vectors
of norm \(a\ne0\) is
\[
                 q^{2m-1}-\eta_{\mathfrak p}^{\,m}q^{m-1}.
                                                               \tag{MF.10}
\]
In the split case choose a nonzero first component and a second
component with prescribed scalar product.
In the nonsplit case the one-coordinate norm has one preimage of
zero and \(q+1\) preimages of every nonzero value.
Convolution gives by induction the zero count
\(q^{2m-1}+(-1)^m(q-1)q^{m-1}\), and (MF.10).
Every nonzero residue vector is a smooth point of the norm map,
since its derivative is a nonzero trace pairing. Each lifting
step contributes \(q^{2m-1}\) choices.

Counting orthonormal columns gives
\[
 |U_m(\mathbb F_q)|=q^{m^2}\prod_{j=1}^m(1-(-1)^jq^{-j})
\]
in the nonsplit case. In the split case the group is
\(\mathrm{GL}_m(\mathbb F_q)\), with order
\(q^{m^2}\prod_{j=1}^m(1-q^{-j})\).
The correction \(Z+Z^*=Y\) lifts every residue isometry, with
\(q^{m^2}\) choices at each step. Thus, for a unimodular lattice,
\[
                         \alpha_L=A_m(\mathfrak p).   \tag{MF.11}
\]

For a nonzero local integer \(n\), put
\(e=\operatorname{ord}_{\mathfrak p}n\).
The normalized norm-vector density \(\alpha(n)\) is \(e_m\) when
\(e=0\). The primitive zero-residue vectors have density
\[
 1+\eta_{\mathfrak p}^{\,m}(q-1)q^{-m}-q^{1-2m}
        =e_m(1+\eta_{\mathfrak p}^{\,m}q^{1-m}).
\]
Divisible vectors contribute \(q^{2-2m}\alpha(n/\pi^2)\) if \(e\ge2\),
and zero if \(e=1\). Consequently
\[
 \frac{\alpha(n)}{e_m}
       =\sum_{j=0}^{e}
              (\eta_{\mathfrak p}^{\,m}q^{1-m})^j.     \tag{MF.12}
\]
This includes dyadic unramified places. For \(m\ge2\) it is
nonnegative and uniformly bounded, since the geometric ratio has
absolute value at most \(1/2\).
This is the actual \texttt{alphaN\_closed} supplier, stronger than
the later difference estimate used by the source.

\paragraph{MF7: arbitrary integral norm densities are bounded and have mean one.}
For any integral nondegenerate \(H\), put
\[
 \alpha_{L,k}(n)=q^{-k(2m-1)}
       \#\{x\in L/\pi^kL:h(x,x)\equiv n\bmod\pi^k\}.   \tag{MF.13}
\]
For \(m\ge1,\ n\ne0\), this stabilizes at sufficiently large \(k\).
For \(m\ge2\) it is bounded uniformly in \(n,k\), and its limit is
bounded, nonnegative, locally constant away from zero, and satisfies
\[
                         \int_R\alpha_L(n)\,dn=1,     \tag{MF.14}
\]
where \(R\) has volume one.

For a primitive vector \(x\), the norm derivative is
\(z\mapsto\operatorname{Tr}h(x,z)\).
If every coefficient were divisible by \(\pi^{d+1}\), perfectness of
the trace pairing and the adjugate identity would make \(x\)
divisible by \(\pi\). Its image is therefore \(\pi^tR\) with
\(0\le t\le d\). Partition primitive vectors into balls modulo
\(\pi^{d+1}L\). On each ball choose a coordinate with derivative
valuation \(t\). Successive correction in that coordinate gives a
uniform map onto a coset modulo \(\pi^{d+1+t}\); the other
\(2m-1\) coordinates are free. There are finitely many balls,
so their pushforward has a bounded locally constant density \(P(n)\).

Writing a nonzero vector as \(\pi^r x\) with \(x\) primitive gives
\[
 \alpha_L(n)=\sum_{\substack{r\ge0\\n\in\pi^{2r}R}}
       q^{-(2m-2)r}P(n/\pi^{2r}).                    \tag{MF.15}
\]
For nonzero \(n\) this is finite. For \(m\ge2\) it converges uniformly
even at zero. The same bound applies to finite-level densities:
the last ball with \(2r\ge k\) contributes at most
\(q^{k-2m\lceil k/2\rceil}\le1\).
These observations prove convergence, boundedness and local
constancy away from zero. The source's sufficient stabilization
bound \(k\ge2d+2\operatorname{ord}_{\mathfrak p}(n)+2\) follows
from these same balls and their scaled versions.
Every vector modulo \(\pi^k\) has one norm residue, whence
\(q^{-k}\sum_{n\bmod\pi^k}\alpha_{L,k}(n)=1\).
Dominated convergence, equivalently the geometric tail in (MF.15),
proves (MF.14). No interchange over a singular fiber is assumed.

\paragraph{MF8: the finite local fibers and their weighted orbit identity.}
Fix a nonzero vector \(x_0\) of norm \(n\in R\setminus\{0\}\), and
write \(V=Ex_0\perp W\), \(r=\dim_EW=m-1\).
For a lattice \(T\) in \(W\), let \(f_L(T)\) count the lattices
\(\Lambda\) locally isometric to \(L\), containing \(x_0\), with
\(\Lambda\cap W=T\). These are exactly the source \texttt{fibLoc}
sets after MF3. Each fiber is finite, because
\[
 Bx_0\oplus T\ \subseteq\ \Lambda\
                  \subseteq\ n^{-1}Bx_0\oplus T^\#.   \tag{MF.16}
\]
The first inclusion is immediate. For the second, projection onto
\(x_0\) has coefficient \(h(x_0,z)/n\in n^{-1}B\), and the
orthogonal component pairs integrally with \(T\).
Both endpoints in (MF.16) are full lattices, with finite quotient.

There are only finitely many local isometry classes of possible
\(T\). Let \(e=\operatorname{ord}_{\mathfrak p}n\) and
\(\ell=\log_q[\Lambda:Bx_0\oplus T]\), with the quotient read
in the containing direction. Projection identifies this quotient
with a subquotient of \(n^{-1}B/B\), so \(0\le\ell\le2e\).
Determinants give \(e+d_T=d_L+\ell\), whence \(d_T\le d_L+e\).
The finite list of Jordan exponents in MF2 therefore supplies the
required finite \texttt{IsLocReps}. We retain actual lattices in
\(W\), not abstract rational isometry classes.

Let \(G=U(V)\), \(H=U(W)\), with the Haar measures of MF5.
The measure on \(V\) which in an \(L\)-basis is
\(|\det H_L|\,dx\) is independent of that basis, and gives \(L\)
volume \(q^{-d_L}\).
Divide this measure by \(dn\) under the norm map to obtain an
invariant measure \(\sigma_n\) on the smooth fiber \(h(x,x)=n\).
Its value on the vectors in \(L\) is
\[
                      \sigma_n(L)=q^{-d_L}\alpha_L(n).
\]
By MF3 the fiber is a single homogeneous space \(H\backslash G\).
The quotient of the two Haar measures is therefore
\(\kappa\,\sigma_n\) for a positive scalar \(\kappa\) independent
of \(L\). This uniqueness can be proved by integrating a compactly
supported function along \(H\): it identifies invariant integration
on the quotient with Haar integration on \(G\).

A right \(K_L\)-orbit in \(H\backslash G\) has measure
\(\nu(L)/\operatorname{vol}(H\cap gK_Lg^{-1})\).
Regroup these orbits by their cut lattice \(T\). The group \(K_T\)
acts on the finite fiber in (MF.16), so finite orbit--stabilizer
gives
\[
 \sum_{[\Lambda]\text{ over }[T]}
 \frac{1}{\operatorname{vol}(\operatorname{Stab}_H\Lambda)}
                  =\frac{f_L(T)}{\nu(T)}.
\]
Consequently, with the literal definitions
\[
 \lambda_{\mathfrak p}(T)=\frac{A_rq^{r d_T}}{\alpha_T}
                         =\frac{A_r}{\nu(T)},\qquad
 \Psi_{\mathfrak p}(L,n)=\sum_{[T]}f_L(T)\lambda_{\mathfrak p}(T),
\]
we obtain
\[
 \Psi_{\mathfrak p}(L,n)
   =A_r\kappa\,\frac{q^{(m-1)d_L}\alpha_L(n)}{\alpha_L}.
                                                               \tag{MF.17}
\]
This proves \texttt{LIRatio} for any two integral lattices.
It also covers empty representation fibers: both sides are zero.
No cancellation by a possibly zero representation density was used.

\paragraph{MF9: pin the scalar, scale the form, and obtain literal LIStat.}
First suppose \(n=1\). Choose \(L_0=Bx_0\perp T_0\) with
\(d_{L_0}=d_{T_0}\le1\), as in MF2.
Every integral lattice of this local type containing \(x_0\)
splits orthogonally as \(Bx_0\perp T\): the projection coefficient
\(h(x_0,z)\) is integral. Its cut has the same rank and determinant
valuation as \(T_0\), so MF2 makes it integrally isometric to \(T_0\).
For a fixed cut, the lattice is uniquely determined. Hence
\[
                       \Psi_{\mathfrak p}(L_0,1)
                              =\lambda_{\mathfrak p}(T_0).       \tag{MF.18}
\]

The finite isometry group modulo \(\pi^k\), for \(k>d_{L_0}\),
acts transitively on all vectors of norm \(1\) modulo \(\pi^k\).
Such a vector lifts to norm exactly \(1\): its norm derivative is
surjective since its norm is a unit, so correction preserves its
residue modulo \(\pi^k\). Its integral orthogonal complement is
again of the unique minimal type just described.
The stabilizer of the first vector consists exactly of the
isometries of the cut modulo \(\pi^k\); the first row and column
are forced by its norm being one. Thus
\[
 a_{L_0}(k)=n_{L_0}(k,1)a_{T_0}(k),\qquad
                    \alpha_{L_0}=\alpha_{L_0}(1)\alpha_{T_0}.
                                                               \tag{MF.19}
\]
The powers of \(q\) match because
\(m^2=(2m-1)+(m-1)^2\). Positivity of the matrix densities shows
\(\alpha_{L_0}(1)>0\). Equations (MF.17)--(MF.19) therefore give
\(\kappa=1\) at norm one.

For general \(n\), use \(h_0=n^{-1}h\), for which \(x_0\) has norm
one. Write \(e=\operatorname{ord}_{\mathfrak p}n\).
Apply the following scaling comparison on a reference lattice
integral for \(h_0\), constructed below. The originally supplied
\(h\)-integral \(L\) need not be integral for \(h_0\); the
lattice-independent scalar \(\kappa\) in MF8, equivalently
\texttt{LIRatio} with the nonzero reference density, then transfers
the computed constant to that original \(L\).
Multiplication of a rank-\(s\) Gram matrix by \(n\) changes the
determinant valuation by \(se\), and comparison of the congruences
at levels \(k\) and \(k-e\) gives
\[
 \alpha_{n h_0,L}=q^{s^2e}\alpha_{h_0,L},\qquad
 \alpha_{n h_0,L}(n)=q^e\alpha_{h_0,L}(1).             \tag{MF.20}
\]
For the first equality, the \(s^2\) Hermitian equations lose \(e\)
digits; for the second, the single norm equation loses \(e\)
digits. More explicitly the free lifts of matrices contribute
\(q^{2s^2e}\), and those of vectors contribute \(q^{2se}\);
division by their respective normalizing powers gives (MF.20).
The isometry groups, cuts and fibers do not change when the form
is scaled. Also \(\lambda_{\mathfrak p}\) is unchanged, since its
two scaling factors \(q^{s^2e}\) cancel.
The ratio in (MF.17) changes by \(q^{-(m-1)e}\).
Its pinned scalar is thus \(q^{(m-1)e}\).

The locally chosen reference lattice is available in the literal
global formulation too. Approximate a local basis of the
complement by rational vectors, preserving its integral Gram
elementary divisors, and multiply these vectors by an integer of
\(F\) which is a unit at \(\mathfrak p\) and clears all other
pairing denominators. Keep \(x_0\) itself unchanged.
This gives a globally integral \(h_0\)-lattice with the specified
localization, and hence a globally integral \(h=n h_0\)-lattice.
We have proved the exact formula
\[
 \boxed{\displaystyle
 \Psi_{\mathfrak p}(L,n)=
 q^{(m-1)(\operatorname{ord}_{\mathfrak p}n+d_L)}
       A_{m-1}(\mathfrak p)\frac{\alpha_L(n)}{\alpha_L}.}
                                                               \tag{MF.21}
\]
This is every clause of \texttt{LIStat}, with the original
\texttt{VecSplit}, \texttt{fibLoc}, and density definitions.
For \(m=1\) the cut is zero dimensional, its automorphism group,
density and empty product all equal one. The argument still applies.

\paragraph{MF10: the rank-one mass constant is independent of the lattice.}
On a positive Hermitian line write \(h(x,y)=a\bar x y\), \(a\in
F_{>0}\), and identify a full lattice with a fractional ideal \(I\).
The rational unitary group is always
\(E^1=\{z:z\bar z=1\}\), independently of \(a,I\).
At each finite place its local stabilizer is also independent of
\(a,I\): it is the norm-one unit subgroup.
A lattice \(J\) is in the genus of \(I\) exactly when the relative
fractional ideal \(JI^{-1}\) has exponent zero at inert primes and
opposite exponents at a split pair. MF3 identifies this completed
description with the source rational-local description.
Hence the set of genus classes is the same quotient of the group
of these balanced ideals by \(\{z\mathcal O_E:z\bar z=1\}\)
for every \(a,I\). Its automorphism group is always
\(\mathcal O_E^\times\cap E^1\), a finite group by the bounded-image
argument of CR4. Finiteness of the genus and positivity of its mass
are IG9. Thus all positive integral Hermitian lines have the same
positive mass \(c_1\).

For rank one with determinant valuation \(d\), the matrix congruence
is \(N(z)\equiv1\bmod\pi^{k-d}\) when \(k>d\).
Counting the \(2d\) free digits above a norm-one residue gives
\(\alpha_L=q^d(1-\eta_{\mathfrak p}q^{-1})\).
Therefore \(\lambda_{\mathfrak p}(L)=1\) at every place.
We have proved precisely
\[
                  \operatorname{MassConst}(E,F,1,c_1),\qquad c_1>0.
                                                               \tag{MF.22}
\]
Its numerical value is unnecessary: it will cancel in the adjacent
rank ratio.

\paragraph{MF11: the real lattice-point constant, without a mass theorem.}
For a rank-\(m\) integral Hermitian lattice \(L\), let
\(\mathfrak d(L)\) be its determinant ideal and put
\(D(L)=\sqrt{N_E\mathfrak d(L)}>0\).
Hermitian symmetry makes the two split determinant valuations
equal, and an inert prime has norm \(q^2\); hence
\[
                         D(L)=\prod_{\mathfrak p}q^{d_L}.        \tag{MF.23}
\]
Only finitely many factors differ from one.

Choose an orthogonal \(E\)-basis \(b_i\) inside \(L\), with
\(A=\bigoplus_i\mathcal O_E b_i\subset L\) of finite index \(I\).
In its complex coordinates, the region
\(h_\tau(z,z)\le T\) for every \(\tau:F\hookrightarrow\mathbb R\)
is a product of ellipsoids. Its volume is
\[
 T^{fm}\left(\frac{\pi^m}{m!}\right)^f
                 \Big/\prod_\tau\det H_\tau.
\]
The covolume of \(A\) in these coordinates is \(v^{fm}\), and
that of \(L\) is \(v^{fm}/I\).
Local determinants, or the determinant of a finite-index basis
change, give
\[
                  D(L)I=\prod_\tau\det H_\tau.        \tag{MF.24}
\]
This also proves that \(D(L)\) is constant on a genus.
The number of lattice points in an expanding bounded ellipsoid,
divided by its scale to the ambient dimension, tends to volume
divided by covolume. To see this, tile by a fixed fundamental
parallelotope. Tiles intersecting the boundary lie in a bounded-width
boundary neighborhood; after rescaling, its volume tends to zero.
The same argument works for this finite product of ellipsoids.
Consequently
\[
 \lim_{T\to\infty}T^{-fm}
 \#\{x\in L:h_\tau(x,x)\le T\text{ for every }\tau\}
       =\frac{(\pi^m/m!)^f}{v^{fm}D(L)}.              \tag{MF.25}
\]
This is the source \texttt{tendsto\_ballCount}, including its powers
of \(2,\pi,v\). Along \(T=j^2\) it has exactly the source's exponent
\(j^{2fm}\).

\paragraph{MF12: the global weighted representation count is a cut-lattice count.}
For \(n\in\mathcal O_F\) totally positive, define
\[
 R_L(n)=\sum_{[M]\in\operatorname{genus}(L)}
       \frac{\#\{x\in M:h(x,x)=n\}}{|\operatorname{Aut}(M)|}.
                                                               \tag{MF.26}
\]
The sums and all vector sets are finite by IG9 and positivity.
Suppose \(m\ge2\). MF4 supplies a rational \(x_0\) of norm \(n\);
let \(W=x_0^\perp\).
Apply finite orbit--stabilizer to the action of
\(\operatorname{Aut}(M)\) on its norm-\(n\) vectors, then apply
MF3 to move each rational vector to \(x_0\).
One obtains the sum of reciprocal stabilizer orders over lattices
\(\Lambda\) in the genus containing \(x_0\), modulo \(U(W)(F)\).
The cut \(T=\Lambda\cap W\) is a full integral lattice.
For a fixed cut \(T\), its possible extensions \(\Lambda\) lie
between the two global lattices in (MF.16), so form a finite set
\(\mathcal F_T\). The finite group \(\operatorname{Aut}(T)\) acts on
it, and its stabilizer is exactly
\(\operatorname{Aut}(\Lambda,x_0)\): the extension of an isometry
of \(T\) fixes \(x_0\), and conversely.
Thus the second orbit--stabilizer calculation gives
\[
                       R_L(n)=\sum_{[T]}
                    \frac{|\mathcal F_T|}{|\operatorname{Aut}(T)|}.
                                                               \tag{MF.27}
\]
In particular this identity has the reciprocal automorphism weights;
an unweighted count of complements would be wrong.

\paragraph{MF13: patch the actual fibers and regroup by complement genera.}
For a fixed rational cut lattice \(T\),
\[
                    |\mathcal F_T|
                          =\prod_{\mathfrak p}f_L(T_{\mathfrak p}).
                                                               \tag{MF.28}
\]
A global lattice determines its localizations. Conversely the finite
local choices in (MF.16) patch by intersection to one global
\(\mathcal O_E\)-lattice with exactly those localizations.
The finite lower and upper bounds in (MF.16) allow the ordinary
Chinese-remainder proof of this assertion, as in LP1--LP3.
Containing \(x_0\), integrality, and the exact cut are all local
conditions, so are preserved by this patching. Rational local
isometry is recovered by MF3.
Outside a finite set containing the primes dividing \(n\), the
determinant of \(L\), and that of one reference cut, both spaces
are unimodular and \(n\) is a unit. Projection then splits the
lattice, and the fiber is a singleton. Thus (MF.28) is a genuinely
finite product. If a local fiber is empty, both sides are zero.

Assume, as an induction hypothesis, that in rank \(m-1\)
\[
                  \operatorname{mass}(T)
                    =c_{m-1}\prod_{\mathfrak p}\lambda_{\mathfrak p}(T).
                                                               \tag{MF.29}
\]
Group (MF.27) by the genus of \(T\). The local fiber cardinalities
and local \(\lambda\)'s are constant on each such genus.
Every finite tuple of allowed local cut types occurs globally:
choose actual representatives in the fixed \(W_{\mathfrak p}\)
and patch them to a lattice in the same rational space \(W\).
This assertion patches lattices; it does not claim the existence
of a new rational Hermitian space with arbitrary invariants.
The finite sum over choices therefore distributes into a product:
\[
                       R_L(n)
                   =c_{m-1}\prod_{\mathfrak p}\Psi_{\mathfrak p}(L,n).
                                                               \tag{MF.30}
\]
The case where a local choice is impossible is again the zero
identity. This proves \texttt{repMass\_eq\_prod\_psi} from its actual
global fiber sets, not from a presumed representation formula.

\paragraph{MF14: the factor identity and the finite singular series.}
Define the literal local and global mass factors and singular series:
\[
 \lambda_{\mathfrak p}(L)=
       \frac{A_m(\mathfrak p)q^{m d_L}}{\alpha_L},\qquad
 \lambda(L)=\prod_{\mathfrak p}\lambda_{\mathfrak p}(L),\qquad
 \mathfrak S_L(n)=\prod_{\mathfrak p}
                        \frac{\alpha_L(n)}{e_m(\mathfrak p)}.
                                                               \tag{MF.31}
\]
The first product is finite because (MF.11) gives
\(\lambda_{\mathfrak p}=1\) when \(d_L=0\).
For nonzero \(n\), the second is also finite: outside the primes
dividing \(n\mathfrak d(L)\), (MF.12) gives \(\alpha_L(n)=e_m\).

Using \(A_m=A_{m-1}e_m\), the local identity (MF.21) becomes
\[
 \Psi_{\mathfrak p}(L,n)=
 q^{(m-1)\operatorname{ord}_{\mathfrak p}n}
       \lambda_{\mathfrak p}(L)\,q^{-d_L}
                   \frac{\alpha_L(n)}{e_m(\mathfrak p)}.
                                                               \tag{MF.32}
\]
Multiplication, with (MF.23) and unique ideal factorization, gives
\[
 R_L(n)=c_{m-1}\lambda(L)D(L)^{-1}
                  |N_{F/\mathbb Q}n|^{m-1}\mathfrak S_L(n).
                                                               \tag{MF.33}
\]
This is the source \texttt{repMass\_eq}. Every product in this
identity has finite support; no analytic Euler-product interchange
has yet been used.

\paragraph{MF15: weighted ideal boxes and finite-place equidistribution.}
Let
\[
 \mathcal B_T=\{n\in\mathcal O_F:0<\tau(n)\le T
                                  \text{ for every }\tau:F\hookrightarrow\mathbb R\}.
\]
For a nonzero ideal \(\mathfrak a\), a residue class modulo a fixed
ideal coprime to \(\mathfrak a\) is a translate of the product ideal
inside the real Minkowski lattice. Tiling by a fundamental
parallelotope and using Riemann sums gives the weighted asymptotic
\[
 T^{-fm}\sum_{\substack{n\in\mathcal B_T\cap\mathfrak a\\
                         n\equiv b\bmod\mathfrak c}}
       |Nn|^{m-1}
       \longrightarrow
       \frac{m^{-f}}{\sqrt{\operatorname{disc}F}\,
                                  N\mathfrak a\,N\mathfrak c}.
                                                               \tag{MF.34}
\]
Indeed the weight rescales to
\(\prod_\tau x_\tau^{m-1}\) on \((0,1]^f\), whose integral is
\(m^{-f}\); the lattice covolume is
\(\sqrt{\operatorname{disc}F}\,N\mathfrak a\,N\mathfrak c\).
The boundary has measure zero, so no uniform error for varying
ideals is needed for this fixed-ideal limit.

We do need a uniform domination. Distinct points of \(\mathfrak a\)
cannot lie in the same half-open cube of side
\((N\mathfrak a)^{1/f}\): their nonzero difference would have norm
strictly smaller than \(N\mathfrak a\), although its principal ideal
is contained in \(\mathfrak a\).
If \(\mathcal B_T\cap\mathfrak a\ne\varnothing\), then
\(N\mathfrak a\le T^f\). Subdividing the box therefore proves
\[
 \#(\mathcal B_T\cap\mathfrak a)
      \le\frac{2^fT^f}{N\mathfrak a},\qquad
 T^{-fm}\sum_{n\in\mathcal B_T\cap\mathfrak a}|Nn|^{m-1}
      \le\frac{2^f}{N\mathfrak a}.                    \tag{MF.35}
\]

Fix a finite set \(S\) of places and put
\(g(n)=\prod_{\mathfrak p\in S}\alpha_{L,\mathfrak p}(n)/
e_m(\mathfrak p)\). For \(m\ge2\), MF7 makes \(g\) bounded.
Away from \(n\in\mathfrak p^j\), each factor is constant on
sufficiently fine residue classes. Replacing it there by the
corresponding finite step function introduces an error bounded
by a constant times \(1_{n\in\mathfrak p^j}\).
For \((\mathfrak a,\mathfrak p)=1\), (MF.35) bounds its normalized
weighted sum over \(\mathfrak a\) by
\(C/(N\mathfrak a\,q^j)\).
Apply (MF.34), the Chinese remainder theorem, and then let \(j\)
increase. The mean-one identity (MF.14) gives
\[
 T^{-fm}\sum_{n\in\mathcal B_T\cap\mathfrak a}
        |Nn|^{m-1}g(n)
 \longrightarrow
 \frac{m^{-f}}{\sqrt{\operatorname{disc}F}\,N\mathfrak a}
                   \prod_{\mathfrak p\in S}e_m(\mathfrak p)^{-1}
                                                               \tag{MF.36}
\]
whenever \(\mathfrak a\) is prime to \(S\).
The same sums are bounded in absolute value by \(C/N\mathfrak a\),
uniformly in \(T\) and \(\mathfrak a\). This proves the actual
finite-place weighted averaging statement, including the treatment
of the exceptional zero residue.

\paragraph{MF16: the singular-series mean by an absolutely dominated divisor sum.}
Choose \(S\) containing the primes with \(d_L\ne0\).
For ideals prime to \(S\), define
\(\chi_m(\mathfrak a)=
 \prod_{\mathfrak p\notin S}\eta_{\mathfrak p}^{\,m\operatorname{ord}_{\mathfrak p}\mathfrak a}\).
The exact good-prime formula (MF.12) gives, for \(n\ne0\),
\[
 \mathfrak S_L(n)=g(n)
      \sum_{\substack{\mathfrak a\mid(n)\\(\mathfrak a,S)=1}}
                \frac{\chi_m(\mathfrak a)}{(N\mathfrak a)^{m-1}}.
                                                               \tag{MF.37}
\]
The normalized sum of a fixed divisor's contribution is bounded
by \(C/(N\mathfrak a)^m\), using (MF.35).
This is summable for \(m\ge2\): there are at most \(f\) prime ideals
over a rational prime, each of norm at least that prime, so
\(\sum_{\mathfrak p}(N\mathfrak p)^{-m}<\infty\).
Expanding the resulting absolutely convergent Euler product proves
\(\sum_{\mathfrak a}(N\mathfrak a)^{-m}<\infty\).
We may therefore pass the limit through the divisor sum.
Using (MF.36) and then its absolutely convergent Euler product gives
\[
 \begin{split}
 &\lim_{T\to\infty}T^{-fm}
       \sum_{n\in\mathcal B_T}|Nn|^{m-1}\mathfrak S_L(n)\\
 &\hspace{12mm}=
 \frac{m^{-f}}{\sqrt{\operatorname{disc}F}}\,
       \prod_{\mathfrak p\in S}e_m(\mathfrak p)^{-1}
       \sum_{(\mathfrak a,S)=1}\frac{\chi_m(\mathfrak a)}{(N\mathfrak a)^m}
   =\frac{m^{-f}}{\sqrt{(2v)^f}}\,L_m.                \tag{MF.38}
 \end{split}
\]
This proves \texttt{tendsto\_singS\_sum}, with the source's actual
constant. The infinite product was not interchanged with an
uncontrolled average: (MF.35) and the summable divisor majorant
justify the passage, including at the induction's endpoint \(m=2\).

\paragraph{MF17: compare two counts of the same vectors.}
Let \(s\) be one finite set of genus representatives for \(L\).
A nonzero vector in a totally positive definite space has a norm
in \(\mathcal B_T\) exactly when all its archimedean norms are
at most \(T\); the zero vector is the only vector of norm zero.
Hence the following identity holds before taking any limit:
\[
 \sum_{M\in s}\frac{\#\{x\in M:h_\tau(x,x)\le T\ \forall\tau\}}
                       {|\operatorname{Aut}(M)|}
              =\operatorname{mass}(L)+\sum_{n\in\mathcal B_T}R_L(n).
                                                               \tag{MF.39}
\]
The left side, divided by \(T^{fm}\), tends by MF11 to
\[
            \operatorname{mass}(L)
                        \frac{(\pi^m/m!)^f}{v^{fm}D(L)}.
\]
Under the rank-\((m-1)\) induction hypothesis, MF14 and MF16 make
the right side tend to
\[
            c_{m-1}\lambda(L)D(L)^{-1}
                             \frac{m^{-f}}{\sqrt{(2v)^f}}L_m.
\]
The separate zero-vector mass tends to zero because \(fm>0\).
All quantities cancelled below are positive: \(D(L),v,\pi,m!\)
are positive, as is \(c_{m-1}\).
Equality of the two limits and \(m!=m(m-1)!\) therefore prove
\[
 \operatorname{mass}(L)=c_{m-1}R_m\lambda(L),\qquad
 R_m=\frac{\big((m-1)!(v/\pi)^m\big)^f}{\sqrt{(2v)^f}}\,L_m .
                                                               \tag{MF.40}
\]
This is \texttt{massConst\_succ}.
The factor \(m^{-f}\) cancels the extra \(m^f\) in the ball
factorial. The discriminant contribution remains exactly
\(\sqrt{(2v)^f}\), with no unidentified measure constant.

\paragraph{MF18: the genus mass formula by induction.}
Start with the positive constant \(c_1\) in MF10 and define
\[
                   c_m=c_1\prod_{j=2}^m R_j
                   \quad(m\ge1).                    \tag{MF.41}
\]
The empty product is one. MF1 gives \(R_j>0\) for all \(j\ge2\),
so each \(c_m>0\).
Induct on \(m\), using MF17 for the successor step.
For every totally positive definite rank-\(m\) Hermitian space
over \(E\), and every full integral lattice \(L\) in it, we obtain
\[
                    \operatorname{mass}(L)=c_m\lambda(L).
                                                               \tag{MF.42}
\]
This is the source \texttt{massConst\_all}.
It is uniform over the rational Hermitian space and over
projective lattices. The local type at bad places is carried by
\(\lambda(L)\); it has not been silently replaced by the
unimodular type. The proof uses only the preceding local counts,
finite genus sums, lattice patching, and the global representation
theorem MF4. In particular it is not an invocation of the desired
Siegel mass formula under another name.

\paragraph{MF19: the literal standard and model mass-ratio endpoints.}
For an \texttt{IsOrthoLattice} of rank \(m\), the supplied global
orthonormal basis identifies the lattice with
\(\mathcal O_E^m\) and the form with the identity matrix.
Thus \(d_L=0\) at every place and (MF.11) gives
\(\lambda(L)=1\). Its mass is \(c_m\).
For two such lattices of ranks \(M\) and \(M-1\), positivity allows
division of (MF.41), giving
\[
 \boxed{\displaystyle
 \frac{\operatorname{mass}(L_M)}{\operatorname{mass}(L_{M-1})}
    =\frac{\big((M-1)!(v/\pi)^M\big)^f}{\sqrt{(2v)^f}}\,
          \operatorname{lValEta}(E,M).}              \tag{MF.43}
\]
We need only \(M\ge3\), exactly the source quantifier.
The argument internally also treats the step from rank one to
rank two. No formula for the absolute rank-one constant is required.

This proves the literal \texttt{SiegelMassRatio} for arbitrary
supplied orthonormal models. Taking the standard coordinate form
and lattice gives \texttt{StandardLattice.SiegelMassRatioCM},
hence the mathematical content of
\texttt{massRatio\_all\_unramified\_cmFields}.
Conversely every orthonormal model is carried to the standard one
by its displayed basis; transporting lattice classes and stabilizers
preserves mass. The model and standard interfaces therefore agree
with exactly the source normalization. All these global assertions
are consequences of the explicitly retained classical
Hasse--Minkowski theorem in MF4.

\paragraph{MF20: the precise consequence for the quantitative window.}
Insert (MF.43) into QW3. For \(k\ge2\) and \(r\ge k\),
\[
        \operatorname{mass}(L_r)=\operatorname{mass}(L_k)
                         \prod_{M=k+1}^{r}R_M.        \tag{MF.44}
\]
The common positive rank-one constant has already cancelled.
Thus QW9's complete quantitative \texttt{window\_lower} now uses
the proved mass ratio rather than a separate unexplained mass-law
premise. Its original scale order, \(0<\epsilon<1/52\), standard
constants, same-\texttt{PreCh} arithmetic data, and the exact
\(d_u,D_s\) hypotheses remain unchanged.
In particular QW6's constant
\(30+(2+2C_0)/\epsilon\) and QW9's choice \(C_0=5\) require no
alteration.

There are three distinct boundaries. First, Hasse--Minkowski is
a named established theorem used here, whose proof has not been
reconstructed. Second, the source's geometric
\texttt{VolumeEqMass2}, \texttt{VolTie}, indefinite volume
normalization, and compatible \texttt{Core}/arithmetic data are
not consequences merely of (MF.43). Third, this section does not
supply the unrelated short-interval prime estimate or finish any
of the RH/GRH global branches. The completed local density law and
global mass ratio must not be relabeled as any of those stronger
statements.



